\chapter{Kesimpulan dan Saran}

\section{Kesimpulan}

Berdasarkan hasil rancang bangun dan pengujian, kesimpulan penelitian adalah
sebagai berikut.
\begin{enumerate}
    \item Purwarupa Rupiah Digital ritel dibangun pada Hyperledger Fabric 2.5
    dengan lima organisasi, \textit{backend} Go, CouchDB, dan PostgreSQL. Model
    distribusi dua jenjang memisahkan Bank Indonesia, bank validator, PJP, dan
    pengguna ritel. Dompet institusional berada di luar tingkatan KYC ritel.
    \item Kebijakan KYC diterapkan sebagai penurunan tingkatan yang
    deterministik. Pelanggan memberikan data identitas; bank/PJP melakukan uji
    tuntas dan keputusan KYC; Bank Indonesia serta pengawas membaca jangkar
    kepatuhan dan peristiwa pengawasan. Data identitas mentah tidak disimpan
    pada buku besar. Risiko tinggi memerlukan EDD dan persetujuan senior.
    \item Tingkatan BASIC, STANDARD, dan MERCHANT menegakkan batas saldo
    maksimum, per transaksi, keluar harian, keluar bulanan, dan masuk bulanan.
    Jalur transfer yang sama menangani transaksi antarpelanggan dan
    pelanggan-ke-pedagang, sehingga pemeriksaan KYC dan limit tidak terduplikasi.
    \item Pengujian otomatis memverifikasi pemetaan KYC, penolakan risiko
    terlarang, persyaratan risiko tinggi, penghitung limit deterministik,
    pemisahan data privat, dan batas peran. Klaim kinerja hanya boleh dibuat
    dari laporan Caliper versi 2.0 dengan beban transfer yang didokumentasikan;
    hasil beban kerja lama tidak digunakan kembali.
\end{enumerate}

\section{Saran}

\begin{enumerate}
    \item Mengulang protokol Caliper pada beberapa perangkat keras, topologi
    jaringan, dan jumlah \textit{worker} yang lebih luas, lalu melaporkan
    konfigurasi, versi \textit{chaincode}, jumlah sukses/gagal, latensi,
    \textit{throughput}, konflik MVCC, dan pengulangan setiap skenario bersama
    hasilnya.
    \item Memperluas pemisahan profil benchmark fungsional. Pengujian penelitian
    ini telah memisahkan transfer ritel, KYC/\textit{onboarding}, operasi
    moneter, kueri pengawasan, dan kebijakan administratif; penelitian lanjutan
    perlu menambahkan variasi ukuran data, distribusi baca/tulis, dan pengulangan
    statistik agar setiap profil dapat dibandingkan secara lebih kuat.
    \item Menambahkan skenario \textit{hot-wallet contention} sebagai
    eksperimen terpisah dari garis dasar tanpa perebutan kunci. Skenario ini
    tidak digunakan pada pengujian utama karena tingkat konsentrasi pedagang dan
    pola pengulangan transaksi belum didasarkan pada jejak transaksi empiris,
    sehingga dapat menambahkan asumsi yang belum tervalidasi. Pengujian lanjutan
    perlu menetapkan konsentrasi akses dari jejak transaksi yang dianonimkan atau
    distribusi sintetis yang parameternya dilaporkan eksplisit, lalu mengukur
    konflik MVCC awal, keberhasilan setelah \textit{retry}, \textit{throughput}
    akhir, serta latensi persentil ke-95 dan ke-99. Mekanisme \textit{retry}
    harus terbatas, memakai \textit{backoff}, dan menjaga idempotensi agar
    transaksi tidak diterapkan lebih dari satu kali.
    \item Memvalidasi angka limit dan distribusi tingkatan bersama Bank
    Indonesia, bank/PJP, serta ahli APU/PPT. Nilai penelitian adalah parameter
    purwarupa dan tidak boleh diperlakukan sebagai ketentuan CBDC.
    \item Mengkaji remunerasi CBDC (\textit{CBDC rate})---berbunga untuk
    \textit{wholesale} dan tanpa bunga untuk ritel---beserta dampaknya terhadap
    kebijakan moneter dan risiko \textit{disintermediation} sebagai studi
    makroekonomi terpisah, melengkapi purwarupa fungsional ini
    \cite{SyarifuddinBakhtiar2021Monetary, Syarifuddin2024OPTIMAL}.
    \item Mengintegrasikan penyedia eKYC nyata dan menerapkan kontrol retensi,
    enkripsi, penghapusan, serta audit akses untuk data identitas di PostgreSQL.
    \item Menguji \textit{private data collection} atau teknik kriptografi
    penjaga privasi untuk mengurangi keterlihatan metadata transaksi tanpa
    menghilangkan kemampuan pengawasan \cite{Pocher2022Privacy,
    Soana2024Central, Said2026Designing}.
    \item Mengembangkan integrasi, interoperabilitas, dan interkoneksi dengan
    infrastruktur pasar keuangan pada penelitian terpisah setelah kebijakan
    buku besar ritel tervalidasi \cite{bi_whitepaper2022}.
    \item Mengembangkan pembayaran tanpa konektivitas sebagai lingkup terpisah
    yang mencakup perangkat aman, pencegahan \textit{double-spending}, dan
    rekonsiliasi sesuai Project Polaris \cite{polaris_security,
    polaris_handbook}.
\end{enumerate}
